\begin{lemma}
\label{lemma-lift-complex}
Let $R$ be a ring. Let $I \subset R$ be an ideal. Let $\mathcal{P}$
be a class of $R$-modules.
Let $K \in D(R)$ and let $E^\bullet$ be a complex of $R/I$-modules
representing $K \otimes_R^\mathbf{L} R/I$. Assume
\begin{enumerate}
\item each $P \in \mathcal{P}$ is a projective $R$-module,
\item $P_1 \in \mathcal{P}$ and $P_1 \oplus P_2 \in \mathcal{P}$
if and only if $P_1, P_2 \in \mathcal{P}$,
\item if $f : P_1 \to P_2$, $P_1, P_2 \in \mathcal{P}$ is surjective
modulo $I$, then $f$ is surjective,
\item $E^\bullet$ is bounded above and $E^i$ is of the form $P/IP$
for $P \in \mathcal{P}$, and
\item $K$ can be represented by a bounded above complex
whose terms are in $\mathcal{P}$.
\end{enumerate}
Then there exists a bounded above complex $P^\bullet$ whose terms
are in $\mathcal{P}$ with $P^\bullet/IP^\bullet$ isomorphic to
$E^\bullet$ and representing $K$ in $D(R)$.
\end{lemma}

\begin{proof}
By assumption (5) we can represent $K$ by a bounded above complex 
$K^\bullet$ whose terms are in $\mathcal{P}$. Then $K \otimes_R^\mathbf{L} R/I$
is represented by $K^\bullet/IK^\bullet$. Since $E^\bullet$ is
a bounded above complex of projective $R/I$-modules by (4), 
we can choose a quasi-isomorphism $\delta : E^\bullet \to K^\bullet/IK^\bullet$
(Derived Categories, Lemma
\ref{derived-lemma-morphisms-from-projective-complex}).
Let $C^\bullet$ be cone on $\delta$
(Derived Categories, Definition \ref{derived-definition-cone}).
The module $C^i$ is the direct sum $K^i/IK^i \oplus E^{i + 1}$
hence is of the form $P/IP$ for some $P \in \mathcal{P}$
as (2) says in particular that $\mathcal{P}$ is preserved under taking sums.
Since $C^\bullet$ is acyclic, we can apply
Lemma \ref{lemma-lift-acyclic-complex} and find a acyclic
lift $A^\bullet$ of $C^\bullet$. The complex $A^\bullet$ is
bounded above and has terms in $\mathcal{P}$. In
$$
\xymatrix{
K^\bullet \ar@{..>}[r] \ar[d] & A^\bullet \ar[d] \\
K^\bullet/IK^\bullet \ar[r] & C^\bullet \ar[r] & E^\bullet[1]
}
$$
we can find the dotted arrow making the diagram commute
by Derived Categories, Lemma
\ref{derived-lemma-morphisms-lift-projective}.
We will show below that it follows from (1), (2), (3)
that $K^i \to A^i$ is the inclusion of a direct summand
for every $i$. By property (2) we see that $P^i = \Coker(K^i \to A^i)$
is in $\mathcal{P}$. Thus we can take
$P^\bullet = \Coker(K^\bullet \to A^\bullet)[-1]$ to conclude.

\medskip\noindent
To finish the proof we have to show the following: Let $f : P_1 \to P_2$,
$P_1, P_2 \in \mathcal{P}$ and $P_1/IP_1 \to P_2/IP_2$ is split
injective with cokernel of the form $P_3/IP_3$ for some
$P_3 \in \mathcal{P}$, then $f$ is split injective.
Write $E_i = P_i/IP_i$. Then $E_2 = E_1 \oplus E_3$.
Since $P_2$ is projective we can choose a map $g : P_2 \to P_3$
lifting the map $E_2 \to E_3$. By condition (3) the map $g$
is surjective, hence split as $P_3$ is projective. Set $P_1' = \Ker(g)$
and choose a splitting $P_2 = P'_1 \oplus P_3$. Then $P'_1 \in \mathcal{P}$
by (2). We do not know that
$g \circ f = 0$, but we can consider the map
$$
P_1 \xrightarrow{f} P_2 \xrightarrow{projection} P'_1
$$
The composition modulo $I$ is an isomorphism. Since $P'_1$ is
projective we can split $P_1 = T \oplus P'_1$. If $T = 0$, then
we are done, because then $P_2 \to P'_1$ is a splitting of $f$.
We see that $T \in \mathcal{P}$ by (2).
Calculating modulo $I$ we see that $T/IT = 0$.
Since $0 \in \mathcal{P}$ (as the summand of any $P$ in $\mathcal{P}$)
we see the map $0 \to T$ is surjective and we conclude that $T = 0$
as desired.
\end{proof}
